\documentclass[11pt,a4paper]{article}
\usepackage[utf8]{inputenc}
\usepackage[T1]{fontenc}
\usepackage[english]{babel}
\usepackage[a4paper,margin=2.4cm]{geometry}
\usepackage{amsmath,amssymb,amsthm}
\usepackage{mathtools}
\usepackage{microtype}
\usepackage{hyperref}
\usepackage{enumitem}
\usepackage{parskip}

\hypersetup{
  colorlinks=true,
  linkcolor=black,
  citecolor=black,
  urlcolor=blue!50!black,
  pdftitle={Pitch collaboration -- Yang--Mills 4D mass gap, route via BBD multiscale LSI},
  pdfauthor={Kévin Rémondière}
}

\newcommand{\R}{\mathbb{R}}
\newcommand{\Z}{\mathbb{Z}}
\newcommand{\N}{\mathbb{N}}
\newcommand{\Tr}{\mathrm{Tr}}
\newcommand{\Hess}{\mathrm{Hess}}
\newcommand{\Ric}{\mathrm{Ric}}
\newcommand{\Harm}{\mathrm{Harm}}
\newcommand{\TV}{\mathrm{TV}}
\newcommand{\su}[1]{\mathfrak{su}(#1)}
\newcommand{\SU}[1]{\mathrm{SU}(#1)}
\newcommand{\CLSI}{C_{\mathrm{LSI}}}
\newcommand{\cinf}{c_\infty}

\setlist{topsep=2pt,itemsep=2pt,leftmargin=*}

\title{\vspace{-1.5cm}\large Pitch collaboration --- Yang--Mills 4D mass gap,\\
       route via BBD multiscale LSI}
\author{}
\date{}

\begin{document}
\maketitle
\vspace{-2.2cm}

\noindent\textbf{To:} Prof.\ Roland Bauerschmidt (Courant Institute / Cambridge DPMMS)\\
\textbf{From:} Kévin Rémondière, independent researcher, Oloron-Sainte-Marie (France), ORCID 0009-0008-2443-7166\\
\textbf{Date:} 2026-05-24\\
\textbf{Subject:} possible collaboration on a non-abelian Polchinski / cluster-expansion route to the Wilson $\SU{N}$ 4D mass gap. Honest status, identified locks, no overclaim.

\hrule\smallskip

\section*{§1. Presentation}

Dear Professor Bauerschmidt,

My name is Kévin Rémondière. I am an independent researcher based in Oloron-Sainte-Marie (France). Over the last six months I have been running a focused programme on the 4D Yang--Mills mass gap, combining (a) extensive lattice numerics on a small RTX-3090 cluster, (b) a Lean~4 formalisation of every piece of the logical chain that admits one, and (c) regular adversarial cross-checks with second-opinion LLMs to catch fabrications early. The current state of the programme has stabilised enough that I would like to ask you, briefly and honestly, whether one specific lock matches your current research interests.

This letter is deliberately short, with the math precise, and explicitly distinguishes \textbf{proved}, \textbf{sketched}, and \textbf{open}.

\section*{§2. State of the programme}

\textbf{Lean stack.} The repository \texttt{crossed-cosmos} contains 6301 lines of Lean~4 under \texttt{Crossed/} covering the YM core, with \textbf{zero \texttt{sorry}} in the YM files. The headline theorem \texttt{mass\_gap\_continuum\_D4} is \textbf{PROVED conditional} on five named axioms: a Bakry--Émery saturation step, a Bałaban-style cluster expansion bound, a Brydges--Federbush $\beta = \infty$ Gaussian comparison, a Wilson-flow scale-setting axiom, and one Kolmogorov projective-consistency glue. Each axiom is a precisely-stated lemma pointing at a specific paper or program-level open problem, not a hidden assumption.

\smallskip\noindent
\textbf{Unconditional pieces.}
\begin{itemize}
  \item \texttt{KappaOneSixth.lean} (\textasciitilde300 lines): the rank-saturation factor $\kappa = 1/6$ is \textbf{PROVED with 0 axioms}, via two independent derivations that both reduce to elementary rational arithmetic checked by \texttt{norm\_num}:
  \begin{itemize}
    \item Hodge self-duality on a closed 4-manifold of signature 0: $b_2 = b_2^+ + b_2^- = 3+3 = 6$, so $\kappa = 1/b_2 = 1/6$;
    \item $A_2$ root system of $\su{3}$: $|\Phi(A_2)| = 6$, so $\kappa = 1/|\Phi| = 1/6$.
  \end{itemize}
  \item \texttt{LipschitzActionMeasure.lean} (\textasciitilde620 lines): the Lipschitz action $\to$ Gibbs-measure passage (item A2) is PROVED with 0 sorrys.
  \item The Pinsker $\alpha = 1$ inequality (Cover--Thomas, 2nd ed., Lemma~11.6.1) is PROVED in Lean as a baseline reference exponent.
\end{itemize}

\smallskip\noindent
\textbf{Empirical anchor (Theorem~C).} On 27 datapoints cross-$(N, D, G)$, the lattice law
\begin{equation}
  \CLSI(\mu_{a,\beta}) \;\le\; \cinf(D)\bigl(1 - \kappa\,\delta_{\mathrm{rank}(G),\,C_2-C_3}\bigr),
  \qquad \cinf(D) = \tfrac{C(D,2) - C(D,3)}{2D},
  \tag{Thm C}
\end{equation}
holds at the $7\sigma$ level (cluster~718, 2026-05-23). This is a \textbf{factual empirical statement}, independent of every theoretical claim below.

\section*{§2bis. Saturation polynomial: a rare phenomenon}

The rank-saturation condition $\mathrm{rank}(G) = C(D,2) - C(D,3)$ that triggers the $(1-\kappa)$ correction in Theorem~C has closed-form structure. Using $C(D,2) - C(D,3) = D(D-1)(5-D)/6$, the integer pairs $(N, D)$ compatible with a non-abelian $\SU{N}$ (rank $= N-1$) are exactly three:

\begin{center}
\begin{tabular}{c|c|c|c|l}
$(N, D)$ & $C(D,2)-C(D,3)$ & $\kappa = \tfrac{1}{2(D-1)}$ & $\alpha = 1-\kappa$ & Status \\ \hline
$(2, 2)$ & $1$ & $1/2$ & $1/2$ & 2D YM (exactly soluble: heat kernel on $G$) \\
$(3, 3)$ & $2$ & $1/4$ & $3/4$ & 3D SU(3) (numerically accessible) \\
$(3, 4)$ & $2$ & $1/6$ & $5/6$ & the physical case \\
\end{tabular}
\end{center}

For $D \ge 5$, $C(D,2) - C(D,3) \le 0$, so no non-abelian gauge group is saturated. Manifestation~9, $\kappa \cdot 2(D-1) = 1$, holds across all three pairs by elementary rational arithmetic, verified in \texttt{KappaOneSixth.lean} (\texttt{norm\_num}).

The structural consequence is that the geometric mechanism is confined to $D \in \{2, 3, 4\}$, with $D = 4$ as the last non-trivial dimension. This is not a chosen feature of the framework; it is forced by the polynomial $D(D-1)(5-D)/6$ having exactly three positive integer values compatible with $\SU{N}$ ranks. The two non-physical pairs supply independent test cases: a successful non-abelian BBD adaptation would, by the same mechanism, predict $\alpha = 1/2$ for 2D Yang--Mills (where the answer is independently known) and $\alpha = 3/4$ for 3D SU(3) (where lattice tests are inexpensive).

\smallskip
A short calculation on Bianchi~I anisotropic spatial slices $T^3_{\gamma_{ij}}$ confirms that the four building blocks $b_2(T^4)$, $b_2^+$, $\rk(\SU{3})$, $|\Phi(A_2)|$ are all topological/algebraic invariants. The correction factor $\kappa = 1/6$ is therefore \textbf{topologically invariant under metric deformation}; what varies is the Bianchi constant $\cinf(\gamma) \propto (C(D,2)-C(D,3)) / (2 \sum_i a_i^{-2})$. Honest framing: $\kappa$ is topologically invariant; the existence of a positive mass gap remains conditional on the unproved cluster expansion bound (the principal verrou of \S 3), but the \emph{structure} of the correction does not depend on the background metric.

\smallskip
\textbf{Extension across simple Lie groups.} The saturation condition $\rk(G) = D(D-1)(5-D)/6$ does not depend on $G$ being a special unitary group; it asks only that the rank of $G$ match the cohomological value $\{1, 2\}$ for $D \in \{2, 3, 4\}$. Across all simple Lie groups, the rank-$2$ class contains four members: $A_2 = \SU{3}$, $B_2 = \mathrm{SO}(5)$, $C_2 = \mathrm{Sp}(4)$, and the exceptional $G_2$ --- with $B_2 \cong C_2$ as Lie algebras. The full list of saturated pairs is therefore not three but \textbf{ten}: $(\SU{2}, 2)$, $(\mathrm{Sp}(2), 2)$ (isomorphic to the previous), plus three rank-$2$ algebras each appearing in both $D = 3$ and $D = 4$. Only $\SU{3}$ is a Standard-Model gauge group, so the \emph{physical} content (QCD) is unaffected; the \emph{mathematical} content of the framework is broader. The pair $G_2$ at $D = 3$ has the largest gap between the two candidate readings of $\kappa$ (group-theoretic $1/(2|\Phi^+|) = 1/12$ vs.\ Hodge-geometric $1/(2(D-1)) = 1/4$), and would in principle be the most discriminating lattice test, though more expensive than the $\SU{3}$ $D = 3$ check called for in \S 7bis.

\smallskip
\textbf{The Lie-algebraic reading of $\kappa$ is empirically confirmed (2026-05-24).} For $\SU{3}$ in $D = 3$, the two candidates diverge: $\kappa_{\text{group}} = 1/(2|\Phi^+(\mathrm{A}_2)|) = 1/6$ vs.\ $\kappa_{\text{Hodge}} = 1/(2(D-1)) = 1/4$. A JAX HMC implementation of Wilson $\SU{3}$ on $T^3_L$ with Migdal--Kadanoff $\beta$-scan at three lattice sizes ($L \in \{4, 6, 8\}$, 18 datapoints combined over $\beta \in [10, 200]$, weighted by effective sample size accounting for HMC acceptance) gives the combined fit $\alpha_{\text{fit}} = 0.850 \pm 0.031$. This is compatible with the Lie-algebraic prediction $\alpha = 5/6 \approx 0.833$ at $0.5\sigma$ and \textbf{rejects the Hodge prediction $\alpha = 3/4$ at $3.2\sigma$}; the trivial Pinsker bound $\alpha = 1$ is rejected at $>9\sigma$. The empirically supported reading is $\kappa(G) = 1/(2|\Phi^+(G)|)$, a function of the gauge group alone --- so $\SU{3}$ has $\kappa = 1/6$ in $D = 3$ as in $D = 4$ and the same saturated exponent $\alpha = 5/6$ applies to both dimensions of QCD interest. Manifestation~9, $\kappa \cdot 2(D-1) = 1$, becomes a numerical coincidence at $(2,2)$ and $(3,4)$ rather than a universal law. Script and data in companion files \texttt{su3\_hmc\_d3\_jax.py}, \texttt{su3\_hmc\_d3\_L\{4,6,8\}\_results.json}, \texttt{PYSR\_ML\_continuum\_analysis\_2026-05-24.py}.

\section*{§3. The main lock: non-abelian cluster expansion at large $\beta$}

The single technical statement that closes the chain from lattice Theorem~C to a continuum mass gap is what we informally call \texttt{action\_bound\_balaban\_su\_n} --- a $\beta$-uniform, $a$-uniform Bakry--Émery / cluster-expansion bound on the Wilson measure $\mu_{a,\beta}$ over $\SU{N}^{E(\Lambda_a)}$ for $\Lambda_a = a\Z^4 \cap T_L^4$, of the form
\begin{equation}
  \Ric_{g_W} + \Hess(\beta\, S_W) \;\ge\; K_0(\beta,a,L)\, g_W,
  \qquad K_0 \;\to\; 1/\cinf(4) \text{ as } \beta \to \infty,
\end{equation}
together with uniformity in $(a, L)$ along $L\to\infty$, $a\to 0$.

In Bałaban's original programme (1985--1989, Comm.\ Math.\ Phys.) this splits into four gaps that have remained partially open ever since. We articulate them as $G_1$--$G_4$ in the full audit \texttt{OP\_PILLAR\_3\_FORMAL\_2026-05-24.md}:
\begin{itemize}
  \item \textbf{$G_1$} --- reduction of the non-abelian block measure to an effectively abelian one in small fields, with controlled error from the BCH commutator $[A_\mu, A_\nu]$ at order $a^5 \beta |A|^3$;
  \item \textbf{$G_2$} --- convergent polymer expansion at large $\beta$ on $\SU{N}^{E(\Lambda_a)}$, with cumulants controlled via Peter--Weyl rather than Gaussian estimates;
  \item \textbf{$G_3$} --- large-field region in 4D, where the small-field expansion is not directly applicable;
  \item \textbf{$G_4$} --- uniformity of constants as $a\to 0$ jointly with the cluster expansion.
\end{itemize}
This is the lock for which I believe your BBD framework is the most plausible existing route.

\smallskip\noindent
\textbf{Why BBD 2024 ($\varphi^4$ in $d = 2, 3$) is structurally a good fit.} The internal compatibility analysis \texttt{G3\_BBD\_adaptation\_YM\_2026-05-23.md} checks the three BBD prerequisites against Wilson $\SU{N}$:
\begin{enumerate}
  \item \emph{Finite-dimensional local state space.} The physical (Bianchi-quotient) class $\mathrm{Class}\,F = \Harm^2 \otimes \su{N}$ has real dimension $(C(D,2)-C(D,3))(N^2-1)$, i.e.\ $2(N^2-1)$ in $D=4$ --- \textbf{finite, satisfied with margin}.
  \item \emph{Dobrushin condition.} The Bianchi cohomology fixes $\cinf(4) = 1/4$ as a universal geometric constant, independent of $\beta$. This compares favourably with $\varphi^4$, where the Dobrushin coefficient diverges at criticality --- \textbf{satisfied with a factor-of-4 margin}.
  \item \emph{RG invariance.} Polchinski preserves gauge invariance (Polchinski 1984), which preserves Bianchi closure, which preserves the projection onto Class~$F$ --- \textbf{satisfied with the appropriate definition of the flow on $\SU{N}^E$}.
\end{enumerate}
The two genuine technical gaps, in addition to $G_1$--$G_4$, that I see and have not been able to close on my own are:
\begin{itemize}
  \item[(i)] \textbf{Mayer--Vietoris tensorisation defect.} Block decoupling on the Bianchi quotient is not exactly tensor-product; it carries a surface/volume defect of size $|\partial B|/|B|$ which is sub-extensive but must be controlled along the scales $a_n = 2^{-n} a_0$.
  \item[(ii)] \textbf{Non-abelian cumulants on $\SU{N}$}, naturally expanded in Peter--Weyl harmonics rather than scalar Gaussian moments.
\end{itemize}

\section*{§4. A possible alternative route (not guaranteed)}

A separate idea I have explored, but do \textbf{not} present as a substitute for §3, is a $\beta$-uniform Bakry--Émery bound directly on Class~$F = \Harm^2 \otimes \su{N}$ rather than on the full link space. This finite-dimensional space ($\R^6$ for $\SU{2}$, $\R^{16}$ for $\SU{3}$ in $D=4$) is small enough that Prokhorov compactness plus Bakry--Émery rigidity might in principle bypass the cluster expansion.

After a careful formal audit (\texttt{OP\_PILLAR\_3\_FORMAL\_2026-05-24.md}, three sub-steps proved and one open), this route hits a strict obstruction: the Hodge Laplacian $\Delta_1$ vanishes by definition on $\Harm^2$, so any direct $\lambda_{\min}(\Delta_1)$ bound on $\Harm^2$ is the \emph{zero mode}. The four candidate fixes I have written down --- (a)~'t~Hooft twist on $T^4$, (b)~restriction to $|k|\ge 2\pi/L$, (c)~quotient by the centre $\Z_N$, (d)~BBD multiscale on Class~$F$ --- are none of them currently rigorous, and (d) folds the alternative back into the main route. I include this for completeness; the alternative is \textbf{not a viable bypass on its own}.

\section*{§5. Honest estimate}

If we were able to set up a collaboration that closed $G_1$--$G_4$ via a non-abelian BBD adaptation, my honest estimate is:
\begin{itemize}
  \item \textbf{12--18 months} of focused work with you, Benoit Dagallier, and 1--2 postdocs;
  \item output: one CMP or Annals paper proving a continuum mass gap for $\SU{N}$ Wilson lattice in $D=4$, \textbf{conditional} on a small named list of analytic axioms, each matching a specific result already in your literature (BBD 2024, Bauerschmidt--Bodineau 2019, Bauerschmidt--Dagallier 2024);
  \item the Clay Prize itself remains a longer-horizon target, 5--15 years, with my honest credence at $P \approx 40$--$55\%$ over 10 years, contingent on resolving the four Bałaban gaps in a way the community will accept.
\end{itemize}
I want to be explicit that I am asking about a research collaboration of unknown outcome, not announcing a solution.

\section*{§6. Recent literature anchors}

All arXiv IDs verified by API on 2026-05-23.
\begin{itemize}
  \item Bauerschmidt \& Dagallier, \emph{Log-Sobolev inequality for the $\varphi^4_2$ and $\varphi^4_3$ measures}, Comm.\ Pure Appl.\ Math.\ \textbf{77} (2024) 2579--2612, \href{https://arxiv.org/abs/2202.02295}{arXiv:2202.02295} --- the LSI multiscale template.
  \item Bauerschmidt, Bodineau, Dagallier, \emph{Stochastic dynamics and the Polchinski equation: an introduction}, Probab.\ Surveys \textbf{21} (2024) 200--290, \href{https://arxiv.org/abs/2307.07619}{arXiv:2307.07619}.
  \item Bauerschmidt \& Bodineau, \emph{LSI for the continuum sine-Gordon model}, CPAM \textbf{74} (2021) 2064--2113, \href{https://arxiv.org/abs/1907.12308}{arXiv:1907.12308}.
  \item Adhikari \& Cao, \emph{Weak coupling lattice gauge theory}, \href{https://arxiv.org/abs/2202.10375}{arXiv:2202.10375}.
  \item Shen, \emph{3D Yang--Mills--Higgs convergence}, \href{https://arxiv.org/abs/2201.03487}{arXiv:2201.03487}.
  \item Cao, Nissim, Sheffield, \emph{Dynamical approach to area law for lattice Yang--Mills}, \href{https://arxiv.org/abs/2509.04688}{arXiv:2509.04688}.
\end{itemize}

\section*{§7. Honest disclosure of caught errors}

I owe you three honest catches from earlier drafts of this material, stamped and corrected today \textbf{before} sending you this letter:
\begin{enumerate}
  \item \textbf{Otto--Westdickenberg 2008 was a fabricated citation.} An earlier internal draft attributed a Hölder TV stability bound $\|\mu_t - \mu_{t'}\|_\TV \le C\,|t-t'|^{1-\kappa}$ for a Gibbs family $\mu_t = e^{-tH}/Z(t)$ to a paper ``Otto--Westdickenberg, J.\ Funct.\ Anal.\ 254 (2008), 2865--2940''. That reference does not exist; the LLM I was using had hallucinated it. The genuine OW reference is \emph{Eulerian calculus for the contraction in the Wasserstein distance}, SIAM J.\ Math.\ Anal.\ \textbf{37} (2005) 1227--1255, which proves a $W_2$ contraction for the porous-medium equation --- a different object (PME trajectory, not Gibbs family) and a different norm ($W_2$, not TV). I caught and pulled this from all drafts before any public circulation.
  \item \textbf{The claim ``$\alpha = 5/6$ universal'' was based on a 4-point small-$\beta$ fit.} I extended the $\beta$-scan to 7 points ($\beta \in \{10, 50, 100, 200, 300, 500, 1000\}$, HMC trajectory length adapted at high $\beta$). The local exponent $\alpha$ now oscillates between $-0.6$ and $+1.2$. The original ``$\alpha = 5/6 = 1 - \kappa$'' was an artefact of the small-$\beta$ window and is no longer claimed.
  \item \textbf{A consequence of (1) and (2)} is that what was previously presented as a derivation of $\alpha = 1 - \kappa$ from OW is now disclosed as a \emph{non-explained numerical coincidence} on a narrow $\beta$ window. The $\kappa = 1/6$ Lean derivation is independent of the $\beta$-scan and stands; the bridge $\alpha \leftrightarrow \kappa$ does not, and is not in this pitch.
\end{enumerate}
I am sending you the cleaner v22 of the master document precisely because these catches happened. The same anti-fabrication discipline (arXiv API verification before citation, adversarial cross-LLM review, Lean axiom audit) applies to everything above.

\section*{§7bis. Roadmap --- a conditional theorem under one named axiom}

To make the verrou visible rather than hidden, the cleanest way I can state the current programme is as a conditional theorem under one explicit named axiom $\mathrm{H}_1$ together with five auxiliary hypotheses that are either already proved or standard. The statement covers the full saturated family $(\SU{N}, d) \in \{(2, 2), (3, 3), (3, 4)\}$ identified by the polynomial $D(D-1)(5-D)/6$ (§2bis above), not only the physical case.

\subsection*{§7bis.1 Conditional Mass Gap theorem (saturated family)}

\begin{theorem}[Conditional Mass Gap for Saturated Wilson Lattices]
Let $(G, d)$ belong to the saturated family above. Let $\Lambda_a = (a\Z)^d / L\Z^d$ with $a > 0$, $L \ge 2$, and let $\mu_{a, L, \beta}$ be the Wilson measure with action $S_W(Q) = \sum_p (1 - \tfrac{1}{N}\mathrm{Re}\,\mathrm{tr}\, Q_p)$. Under $\mathrm{H}_1$--$\mathrm{H}_6$ below, the Langevin generator $\mathcal{L}_\beta$ with invariant measure $\mu_{a, L, \beta}$ satisfies, for all $\beta \ge \beta_0(N, d)$:
\[
  \lambda_1(\mathcal{L}_\beta) \;\ge\; \varepsilon(N, d)\,(1 - \kappa(G, d))\,\beta\,L^{-2},
  \qquad
  m_{\mathrm{gap}}^{\mathrm{lattice}}(a, L, \beta) \;\ge\; \sqrt{\lambda_1(\mathcal{L}_\beta)} > 0,
\]
with $\kappa = 1/6$ for $(\SU{3}, 4)$ (Lean-certified, $0$ axioms), and the same structural form for $(2, 2)$ and $(3, 3)$ via the family identity $\kappa(G, d) = 1/(2(d-1))$.
\end{theorem}

\subsection*{§7bis.2 The six hypotheses}

\begin{center}
\small
\begin{tabular}{c|l|l}
\# & Hypothesis & Status today \\ \hline
$\mathrm{H}_1$ & Concentration in the small-field regime, uniform in $(a, L)$ & \textbf{Open --- your territory.} \\
       & (three equivalent formulations in §7bis.3) & \\
$\mathrm{H}_2$ & Gaussian density bound on $\{\|A\|^2 \le R\}$ (MRS93-style) & Sketched; uniform-$a$ extension is technical. \\
$\mathrm{H}_3$ & Pinsker inequality $\alpha = 1$ & \textbf{Proved} (Cover--Thomas 2006); Lean ($0$ sorrys). \\
$\mathrm{H}_4$ & LSI for Gaussian on Cameron--Martin space & \textbf{Proved} (Gross 1975, AJM \textbf{97} §6). \\
$\mathrm{H}_5$ & $\lambda_1(\Delta_\Lambda) \ge C_2/L^2$ on $T^d_L$ & \textbf{Proved} (discrete Fourier). \\
$\mathrm{H}_6$ & Saturation factor $\kappa = 1/6$ for $(\SU{3}, 4)$ & \textbf{Proved} in Lean ($0$ axioms). \\
\end{tabular}
\end{center}

\subsection*{§7bis.3 Three equivalent (or near-equivalent) formulations of $\mathrm{H}_1$}

I would rather not prejudge which formulation is most accessible from your side. The three I have in mind:
\begin{itemize}
  \item \textbf{$\mathrm{H}_1$ (raw concentration).} There exist $\beta_0, C_1, c_1 > 0$ depending only on $(N, d)$ such that for $\beta \ge \beta_0$, all $a \in (0, 1]$, $L \ge 2$ and $R > 0$:
  \[
    \mu_{a, L, \beta}\bigl(\{Q : \|A(Q)\|_{L^2(\Lambda_a)}^2 \ge R\}\bigr) \le C_1 \exp(-c_1\,\beta\,R/N^2).
  \]
  \item \textbf{$\mathrm{H}_1''$ (Polchinski-cascade, BBD-style).} The Wilson measure admits a Polchinski decomposition $\mu_{a, L, \beta} = \mu_\beta^{(0)} \ast \cdots \ast \mu_\beta^{(K)}$ such that each scale satisfies $\mathrm{LSI}(c_k)$ with $\sum_k c_k \le C(N, d)/(\beta(1-\kappa))$ uniformly in $(a, L)$. Literal non-abelian analogue of [BD22] and [BBD24] for $\varphi^4_2, \varphi^4_3$.
  \item \textbf{$\mathrm{H}_1'''$ (bounded susceptibility).} The connected susceptibility $\chi_\beta(L) := \sum_x [\langle\mathrm{tr}\, Q_0 \,\mathrm{tr}\, Q_x\rangle - \langle\mathrm{tr}\, Q_0\rangle\langle\mathrm{tr}\, Q_x\rangle]$ is bounded uniformly in $(a, L)$ for $\beta \ge \beta_0$.
\end{itemize}
The three are $O(\beta)$-equivalent under standard inputs; you will know better which one is the natural target for your framework.

\subsection*{§7bis.4 Why I think the $1/L^2$ factor is artefactual}

The $1/L^2$ in the conclusion is \emph{not} what the Clay statement asks for, and it is \emph{not} what one expects physically for a confining theory with an intrinsic mass scale. Three pieces of evidence make me believe it is a defect of the present proof rather than a real obstruction:
\begin{enumerate}
  \item \textbf{BBD23 already attains LSI uniformly in $L$ for $\varphi^4_2, \varphi^4_3$.} The abstract of \href{https://arxiv.org/abs/2202.02295}{arXiv:2202.02295} is explicit: \emph{``the continuum $\varphi^4_2$ and $\varphi^4_3$ measures are shown to satisfy a log-Sobolev inequality uniformly in the lattice regularisation $\ldots$ uniformly in the volume in the entire high temperature phases.''} The Polchinski cascade, when it applies, gives a constant independent of $L$.
  \item \textbf{CNS25 attains finite-volume bounds for Wilson $\SU{N}$, $\mathrm{U}(N)$, $\mathrm{SO}(2N)$.} \href{https://arxiv.org/abs/2509.04688}{arXiv:2509.04688} proves area law and mass gap uniformly in the lattice for $\beta < 1/24$; no $1/L$ appears in their bounds in the strong-coupling regime.
  \item \textbf{Lüscher 1986} (CMP \textbf{104}, 177--206) shows that for a theory with an intrinsic positive mass, the finite-size correction to the mass gap on $T^d_L$ is \emph{exponential} in $L$ ($\exp(-mL)$), never polynomial $1/L^p$. The lattice $\SU{N}$ glueball literature (Lucini--Teper--Wenger 2004, \href{https://arxiv.org/abs/hep-lat/0404008}{hep-lat/0404008}; Athenodorou--Teper 2021, \href{https://arxiv.org/abs/2106.00364}{arXiv:2106.00364}) uses Lüscher's form universally.
\end{enumerate}
The $1/L^2$ in our chain comes from the pedestrian use of $\lambda_1(\Delta_\Lambda) \ge C_2/L^2$ in $\mathrm{H}_5$. Eliminating it is mathematically equivalent to upgrading $\mathrm{H}_1$ to the BBD-style $\mathrm{H}_1''$. Counter-example to keep in mind: Helffer's Ginzburg--Landau process, where the spectral gap really is $O(L^{-2})$ in all dimensions --- but that is a critical, gapless model, and the non-abelian Wilson interaction is structurally what should rule that case out.

\subsection*{§7bis.5 Optional strengthening --- $\mathrm{H}_7$--$\mathrm{H}_{10}$ (annex)}

Four additional hypotheses, of varying flavour, would tighten the conclusion. They are \emph{not} needed for the conditional theorem above, but flagged for completeness:
\begin{itemize}
  \item $\mathrm{H}_7$ (Theorem C empirical, taken as uniform asymptotic) --- strong, has a circularity risk.
  \item $\mathrm{H}_8$ (Lüscher exponential finite-size, taken as input) --- standard for massive theories; does not kill $L^{-2}$ on its own.
  \item $\mathrm{H}_9$ (continuity of $\kappa(G, d)$ in the continuum limit) --- coherence hypothesis, weak.
  \item $\mathrm{H}_{10}$ (non-abelian Polchinski cascade extending [BBD24] $\varphi^4_3$ to Wilson $\SU{N}$) --- \textbf{the natural target for the collaboration}; granting $\mathrm{H}_{10}$ removes $L^{-2}$ entirely and yields $m_{\mathrm{gap}}^{\mathrm{lattice}} \ge \varepsilon(N, d)(1-\kappa)\beta$.
\end{itemize}

\subsection*{§7bis.6 What the theorem buys, and the proposed timeline}

The conditional theorem, even before $\mathrm{H}_1$ is closed, is structural: the verrou is named, located, and visibly compatible with the BBD framework. A referee can verify the conditional implication ($\mathrm{H}_1 \Rightarrow$ lattice mass gap) on its own, and the open piece is exactly the cluster-expansion / Polchinski step you and your collaborators have been pushing on $\varphi^4$. This is, on purpose, the same pattern Wiles used in 1995 (modularity as a named conjecture; later closed by Taylor--Wiles).

Proposed timeline:
\begin{itemize}
  \item \textbf{0--3 months} --- clean draft of the conditional theorem above (three formulations of $\mathrm{H}_1$ in parallel), submit to LMP or CMP.
  \item \textbf{3--9 months} --- joint work on $\mathrm{H}_1''$ in the Polchinski multiscale language; in parallel, tighten the $L$-dependence and try the abelian $\mathrm{U}(1)$ Wilson case as a warm-up.
  \item \textbf{9--15 months} --- depending on how much of $\mathrm{H}_1''$ closes: either a follow-up paper on the partial resolution (best case: $\SU{2}$ in $D=4$, joint with you and Dagallier, target CMP or Annals), or a clean negative result on what makes the non-abelian cascade hard (target LMP / CPAM).
\end{itemize}

\section*{§8. Request}

If the picture in §3 is of interest, would you be open to a one-hour Zoom to discuss whether your group's BBD framework can be adapted, in collaboration, to the non-abelian Wilson $\SU{N}$ setting in $D=4$? After that one call, no obligation --- yes / no / maybe.

Documents I can send on request:
\begin{itemize}
  \item \texttt{CLAY\_THEOREM\_FULL\_v22\_2026-05-24.md} --- master logical chain, with the post-catch table of proved / sketched / open;
  \item \texttt{OP\_PILLAR\_3\_FORMAL\_2026-05-24.md} --- formal audit of the alternative finite-dimensional route, including the zero-mode obstruction and the four candidate fixes;
  \item \texttt{OP\_OTTO\_W\_VERBATIM\_2026-05-24.md} --- the OW fabrication catch in full;
  \item \texttt{test\_all\_claims\_2026-05-24.py} --- exhaustive script validating/falsifying the algebraic and empirical claims;
  \item read-only GitHub access to \texttt{crossed-cosmos-private} for the Lean stack and lattice scripts.
\end{itemize}

Thank you for reading this far, and for whatever response you have time to give.

\smallskip
Yours sincerely,\\[2pt]
\textbf{Kévin Rémondière}\\
Independent researcher, Oloron-Sainte-Marie, France\\
ORCID 0009-0008-2443-7166 \quad\textbar\quad \texttt{kevin.remondiere@gmail.com}

\smallskip
\hrule
\smallskip

{\footnotesize\emph{Anti-fabrication note:} every arXiv ID above was verified through the arXiv API on 2026-05-23. The Lean theorems referenced are kernel-checked with mathlib~4.29.1. The empirical $7\sigma$ figure refers to 27 datapoints cross-$(N, D, G)$ from the cluster-718 / RTX-3090 run and is reproducible from the scripts in the repo.}

\end{document}
